\chapter{P1：保证后面一直使用同一版数据}

\section{承上启下}

P0 只证明来源和本地条件可行，尚未形成可复现输入。P1 的业务问题是：\textbf{怎样保证后续每个阶段使用的是同一对评论与元数据文件，而不是某次中断、替换或误发现的文件？}

\begin{beginnerbox}
继续沿用 P0 的运营例子：如果今天的“漏水”结论来自文件 A，下周却不小心读了同名的文件 B，就无法判断变化来自商品还是来自数据。P1 要做的就是给本次输入做一份清单和版本记录；以后要换数据，必须明确换一版。
\end{beginnerbox}

\section{输入、边界与价值}

输入是 P0 核实过的评论和元数据来源角色。P1 不解释评论内容、不选择品类，也不声称拥有服务端官方 checksum。业务价值是让不同时间的分析仍能回答“用了哪批输入”；技术价值是把下载恢复、并发安全、完整性与发布原子性组成一条证明链。这里的“冻结”不是把文件锁死，而是给当前版本盖章，避免下游偷偷换文件。

\section{开发过程}

\subsection{1. 继续下载前先确认还是同一个文件}

断点续传不是简单地从旧文件末尾继续写。流程先比较远端对象身份，例如对象标签或最后修改时间；只有身份仍与已记录对象一致，Range 恢复才有意义。远端身份变化时必须停止，避免把两个版本拼成一个文件。初学者可以先把它理解为：搬货搬到一半回来，先确认仓库里的货没有换批次。

\subsection{2. 保证同一时间只有一个下载者}

同一目标在同一时间只允许一个写入者。独占锁拒绝第二个下载进程；遗留锁不能静默删除，而要经过显式恢复。这里的锁保护的是并发行为，不是数据正确性的替代品。它解决的是“两个人同时往同一个箱子里装东西”的问题。

\subsection{3. 检查文件完整，再留下文件指纹}

下载或人工取得后，先完整读取 gzip 流，检查压缩结构没有截断；再计算本地内容哈希，作为以后重新核对同一字节内容的文件指纹。人工直接下载的记录只支持本地大小、gzip 与哈希均通过，不伪装成已捕获浏览器 HTTP 回执或平台官方校验和。

\subsection{4. 给两类文件做一份输入清单（manifest）}

评论与元数据两个角色全部通过后，才生成组合快照 manifest。manifest 可以先理解为“输入清单”：同时记录角色、来源身份、本地完整性和快照身份。下游读取“reviews”“metadata”等角色映射，而不是扫描原始目录猜哪个文件最新。

\subsection{5. 只有整批检查通过才发布}

manifest 先写临时文件，内容完整并通过检查后再替换正式文件。如果任一输入失败，旧的有效 manifest 不应被半成品覆盖。通俗地说，公开的是一张完整封条，而不是边写边让下游读取的草稿。这类“要么整批成功、要么不换旧版本”的动作，工程上叫原子发布。

\section{本阶段学到的知识}

本阶段把\textbf{可恢复传输、远端对象身份、跨进程锁、gzip 流校验、本地哈希、manifest 和原子发布}串成一个系统。初学者先掌握一句话：先确认没有换批次，再确认文件没损坏，最后把通过检查的整批输入登记下来。单独拥有哈希并不能证明来源官方背书；单独拥有锁也不能证明文件完整。只有这些机制按顺序组合，才能给下游稳定输入。

\begin{beginnerbox}
项目例子：假设原始目录里多出一个名字相似的压缩包。P2 不能因为它“更新”就读取它，而必须从冻结 manifest 的角色字段解析输入。这样文件命名或目录噪声不会悄悄改变数据快照。
\end{beginnerbox}

\section{问题、修复与验证}

问题一是中断下载可能留下可读但不完整的 gzip；修复是完整流校验。问题二是并发写入或遗留锁可能污染目标；修复是默认拒绝第二写入者和显式恢复。问题三是人工下载缺少服务端 checksum；修复是诚实限定为本地观察。最后通过失败场景验证：角色缺失、身份变化或完整性失败时，不发布新快照。

\textbf{小练习：}说明“本地 SHA-256 相同”和“平台提供官方 checksum”为什么不是同一断言。前者证明两次本地字节计算一致，后者还要求可核实的服务端声明。

\section{阶段产出与验收边界}

主要产出为 \filelink{artifacts/manifests/phase-1-download.json} 与 \filelink{docs/reports/phase-1-download.md}。P1 Gate 要求两个最终 gzip 角色完整、本地哈希已登记、组合快照原子发布，并发写入受到约束。它不证明浏览器传输过程、服务端官方 checksum、内容语义质量或公开再分发权。P1 在该边界内为 \statuspassed。

\section{怎样进入下一阶段}

P2 必须按 manifest 中的角色读取输入，并在解析前重算本地哈希。任何未登记文件都不能自动加入审计；若要更换数据，必须重新形成 P1 快照，而不是在 P2 临时换源。

\section{面试怎么讲}

\begin{interviewbox}
“我把下载、登记和发布分开：恢复前检查远端身份，写入用独占锁，完成后验证整个 gzip 并计算本地哈希，最后原子发布同时包含评论和元数据角色的冻结 manifest。下游只按 manifest 角色读取，不扫描原始目录。我不会把本地哈希说成官方 checksum。”
\end{interviewbox}

\textbf{掌握标准：}能够画出“身份检查 $\rightarrow$ 恢复 $\rightarrow$ 锁 $\rightarrow$ gzip $\rightarrow$ 哈希 $\rightarrow$ manifest $\rightarrow$ 原子替换”的失败关闭路径。
